%% ma: L10-B13-0079
%% lop: 10
%% chuong: 5
%% bai: 13
%% dang: 10.13.10
%% loai: TN4
%% muc: TH
%% dap_an: "B"
%% nguon: "chuyen-de-cac-so-dac-trung-cua-mau-so-lieu-khong-ghep-nhom-toan-10-chuong-trinh-moi.pdf (GV. Phan Nhật Linh) (P03), Bài 02 Dạng 2 Phần I Câu 4"
%% kien_thuc: "Lớp 10 Bài 13 (trung vị, tứ phân vị)"
%% trang_thai: chinh-thuc
%% ghi_chu: "Tự duyệt theo yêu cầu giáo viên 22/10/2026."
\begin{ex}
Bảng sau cho biết chiều cao (cm) của một nhóm học sinh:
\begin{center}
\begin{tabular}{|c|c|c|c|c|c|c|c|}\hline
$160$ & $178$ & $150$ & $164$ & $168$ & $176$ & $156$ & $172$\\\hline
\end{tabular}
\end{center}
Các tứ phân vị của mẫu số liệu là
\choice{$Q_1=158;Q_2=164;Q_3=174$}{\True $Q_1=158;Q_2=166;Q_3=174$}{$Q_1=160;Q_2=168;Q_3=176$}{$Q_1=150;Q_2=164;Q_3=178$}
\loigiai{Sắp xếp: $150$; $156$; $160$; $164$; $168$; $172$; $176$; $178$.

$Q_2=\dfrac{164+168}2=166$; $Q_1=\dfrac{156+160}2=158$; $Q_3=\dfrac{172+176}2=174$. Chọn B.}
\end{ex}
